\documentclass[a4paper,10pt]{article}
\usepackage[utf8]{inputenc}
\usepackage{graphicx}
\usepackage{subfigure}
\usepackage{fancyhdr}
\usepackage{listings}
\usepackage{ gensymb }


%opening

\pagestyle{fancy} %eigener Seitenstil
\fancyhf{} %alle Kopf- und Fußzeilenfelder bereinigen
\fancyhead[L]{SEP Abgabegespräch SS2015/16} %Kopfzeile links
\fancyhead[C]{} %zentrierte Kopfzeile
\fancyhead[R]{\includegraphics[width=1cm]{../../template/tugraz_bw}} %Kopfzeile rechts
\renewcommand{\headrulewidth}{0.4pt} %obere Trennlinie
\fancyfoot[C]{} %Seitennummer
\renewcommand{\footrulewidth}{0.4pt} %untere Trennlinie


\lstset{language=C++, basicstyle=\ttfamily\tiny}  

\begin{document}

\section*{Movie Collector}

Entwicklen Sie eine Software die Sie Filme verwalten lässt.

Sollte es zu einer Fehlermeldung bei der Abarbeitung eines Befehles kommen
wird dieser "abgebrochen" und es wird wieder auf eine Eingabe gewartet.
Sollte ein Befehl nicht dem vorgegebenen Format entsprechen wird er einfach ignoriert. Case Sensitive Unterstützung bei den Commandos (nicht den Argumenten) ist optional.
Nachfolgend die zu implementierenden Befehle:

\begin{description}
 \item[add $<$moviename$>$ $<$year$>$ $<$producername$>$] Fügt eine Film der Sammlung hinzu. Der Filmname muss dabei eindeutig sein.
  Sollte ein gleicher Film bereits exisiteren wird "Movie does already exist!" ausgegeben.
 \item[show $[$ $<$moviename$>$]$ $] Zeigt die ganze Filmesammlung (zeilenweise den moviename) wenn nur show eingeben wurde oder den jeweiligen Film (siehe Beispiel) wenn ein moviename angegeben wurde.
 Sollte ein möglicherweise angegebener Film nicht existieren wird "Movie does not exist!" ausgegeben.
 \item[delete $<$moviename$>$ ] Löscht den jeweiligen Film aus der Datenbank. Sollte ein möglicherweise angegebener Film nicht existieren wird "Movie does not exist!" ausgegeben.
\end{description}

Orientieren Sie sich zusätzlich an dem unten vorgegebenen Beispiel. Bei der Eingabe von \texttt{quit} oder EOF soll das Programm beendet werden.

\subsection*{Beispiel}
\begin{verbatim}
> add Simpsons 2010 Steven
> add BatmanBegins 2009 Steve
> show
Simpsons
BatmanBegins
> show Simpsons
Moviename: Simpsons
Year: 2010
Producername: Steven
> show Bat
Movie does not exist!
\end{verbatim}

\subsection*{Aufruf}
\texttt{./movieCollector}

\subsection*{Bewertung}
\begin{itemize}
 \item Dokumentation und Programmierstil: 1 Punkt
 \item Einhalten der OOP Konzepte: 1 Punkt
 \item Struktur \& Korrektheit: 1 Punkt
\end{itemize}

\newpage


\end{document}
